
\section{Resistance Taxonomy}
\label{s:appendix-taxonomy}

This appendix provides the full resistance taxonomy used for constructing resistant-client profiles. The taxonomy includes terminology, resistance dimensions and categories, and detailed subcategory definitions with examples and boundary rules. It synthesizes recurring concepts from prior resistance research rather than reproducing one earlier taxonomy. The literature links are summarized below Table~\ref{tab:resistance-main-categories}; the final definitions, examples, and boundary rules are specific to the questionnaire-based setting used in this thesis.

\subsection{Terminology}

The terms below separate the organization of the taxonomy from the experimental intensity condition.

\noindent\textbf{Resistance.} Client responses that reduce cooperation with the questionnaire-based assessment by making an answer less useful or by disrupting the assessment process.

\noindent\textbf{Dimension.} The part of the dialogue mainly affected by resistance. The \emph{content} dimension concerns the information in the current answer; the \emph{interaction} dimension concerns how the client engages with the assessment process.

\noindent\textbf{Category and subcategory.} A category groups related responses within a dimension: withholding, diverting, reframing, contesting, or controlling. A subcategory is a specific response form within one of these categories. The 12 subcategories are used to construct the resistance profiles.

\noindent\textbf{Intensity.} The experimental setting that mainly controls how often resistance is activated. It also slightly changes expression-level sampling on activated turns.

\subsection{Resistance Dimensions and Categories}

\begin{table}[H]
    \centering
    \tabcap{Resistance dimensions and main categories.}
    \label{tab:resistance-main-categories}
    \renewcommand{\arraystretch}{1.08}
    \begin{tabularx}{\columnwidth}{@{}>{\raggedright\arraybackslash}p{0.18\columnwidth}>{\raggedright\arraybackslash}p{0.28\columnwidth}>{\raggedright\arraybackslash}X@{}}
        \toprule
        \rowcolor{black!8}
        \thead{Dimension} & \thead{Categories} & \thead{Definition} \\
        \midrule
        Content & Withholding;\newline diverting;\newline reframing &
        The response limits usable information, moves away from the questionnaire item, or changes how a profile-supported problem is presented. \\
        Interaction & Contesting;\newline controlling &
        The response challenges, blocks, redirects, or tries to control the questionnaire process. \\
        \bottomrule
    \end{tabularx}
\end{table}

\noindent\textit{Literature links.} Withholding draws on response-quantity resistance, minimum talk, and nonanswer~\cite{otani1989client,chamberlain1984observation,li2026recap}; diverting on verbosity, sidetracking, and inattention~\cite{otani1989client,chamberlain1984observation,li2026recap}; and reframing on minimization, externalization, blaming, and denial~\cite{otani1989client,li2026recap}. Contesting is related to challenging, discounting, and hostility~\cite{chamberlain1984observation,li2026recap,schwartz2019resistance}, while controlling draws on limit setting and pursuit of the client's own agenda~\cite{otani1989client,chamberlain1984observation,li2026recap}.

\FloatBarrier

\subsection{Resistance Subcategories, Examples, and Boundary Rules}

Tables~\ref{tab:content-resistance-subcategories} and~\ref{tab:interaction-resistance-subcategories} use the same three fields. The definition states the main response pattern, the example illustrates how it may appear in a questionnaire answer, and the boundary rule distinguishes it from the closest neighboring subcategory. The examples are illustrative rather than fixed responses used by the patient agent. During generation, the model receives the full profile for the selected subcategory and produces wording for the current question and dialogue context.

\begin{table*}[tb]
    \centering
    \tabcap{Content-level resistance subcategories. Example question: ``Unexplained aches and pains, for example in the head, back, joints, abdomen, or legs?''}
    \label{tab:content-resistance-subcategories}
    \renewcommand{\arraystretch}{1.08}
    \begin{tabularx}{\textwidth}{@{}p{0.20\textwidth}p{0.39\textwidth}X@{}}
        \toprule
        \rowcolor{black!8}
        \thead{Subcategory} & \thead{Definition and example} & \thead{Boundary rule} \\
        \midrule
        \rowcolor{black!5}
        \multicolumn{3}{@{}l}{\textit{Withholding}} \\
        \addlinespace[2pt]
        Vague/minimal information &
        The client gives a short, vague, or unclear answer that does not clearly answer the item. \textit{Example:} ``Maybe a little.'' / ``Hard to say.'' &
        This label fits responses that contain some answer but lack detail or clarity without offering an explanation. It is different from silence/nonresponse, where almost no substantive answer is given. \\
        Silence/nonresponse &
        The client gives no substantive answer to the current item. \textit{Example:} ``...'' / ``No comment.'' / ``Huh...'' &
        The key feature is the absence of a usable answer. If the client explicitly refuses to discuss topic-related symptoms or sets conditions for answering, topic blocking/conditional engagement is more appropriate. \\
        \addlinespace[4pt]
        \rowcolor{black!5}
        \multicolumn{3}{@{}l}{\textit{Diverting}} \\
        \addlinespace[2pt]
        Topic shift/inattention &
        The client moves away from the current question or does not seem to follow what was asked. \textit{Example:} ``Anyway, work has been really stressful lately.'' &
        The response drifts away from the questionnaire item without clearly trying to control the assessment. This differs from agenda redirection, where the client explicitly tries to change what the assessment should discuss. \\
        Irrelevant verbosity &
        The client gives a long response with many irrelevant details, making the answer difficult to identify. \textit{Example:} ``Well, it depends on the chair, the weather, my shoes, what I ate, and whether I slept well...'' &
        The response is long but low in relevance. A long answer should not be categorized this way if it still clearly answers the symptom question. \\
        \addlinespace[4pt]
        \rowcolor{black!5}
        \multicolumn{3}{@{}l}{\textit{Reframing}} \\
        \addlinespace[2pt]
        Minimization/\allowbreak normalization &
        The client acknowledges a problem but presents it as less serious, less important, or normal. \textit{Example:} ``My back hurts sometimes, but everyone gets aches.'' &
        The client acknowledges the symptom but downplays its clinical importance. This differs from denial/dismissal, where the client rejects the presence of the problem. \\
        Externalization/\allowbreak blaming &
        The client explains the problem mainly through external causes or other people. \textit{Example:} ``My body hurts because my boss makes me work too much.'' &
        The main function is to shift the explanation outward. External causes alone are not resistance if the client still clearly describes the symptom and its impact. \\
        Denial/dismissal of profile-relevant problems &
        The client denies or dismisses a problem that should be present according to the clinical profile or earlier dialogue context. \textit{Example:} ``No, I don't have any aches or pains.'' &
        A negative answer is not automatically resistance. This label only applies when the denial conflicts with the clinical profile or previous dialogue context. \\

        \bottomrule
    \end{tabularx}
\end{table*}

\begin{table*}[tb]
    \centering
    \tabcap{Interaction-level resistance subcategories. Example question: ``Unexplained aches and pains, for example in the head, back, joints, abdomen, or legs?''}
    \label{tab:interaction-resistance-subcategories}
    \renewcommand{\arraystretch}{1.08}
    \begin{tabularx}{\textwidth}{@{}p{0.20\textwidth}p{0.39\textwidth}X@{}}
        \toprule
        \rowcolor{black!8}
        \thead{Subcategory} & \thead{Definition and example} & \thead{Boundary rule} \\
        \midrule
        \rowcolor{black!5}
        \multicolumn{3}{@{}l}{\textit{Contesting}} \\
        \addlinespace[2pt]
        Challenging assessment or intent &
        The client questions the reason, relevance, or legitimacy of the question or assessment. \textit{Example:} ``What does this have to do with my mental health?'' &
        The response challenges why the question is being asked. Neutral clarification should not be treated as resistance. \\
        Dismissive/\allowbreak hostile criticism &
        The client criticizes, mocks, or attacks the question, assessment, or assistant. \textit{Example:} ``These questions are so freaking stupid.'' &
        The key feature is dismissive or hostile wording toward the question, scale, assessment, or assistant. This differs from challenging assessment or intent, which may question relevance without hostile wording. \\
        \addlinespace[4pt]
        \rowcolor{black!5}
        \multicolumn{3}{@{}l}{\textit{Controlling}} \\
        \addlinespace[2pt]
        Topic blocking/\allowbreak conditional engagement &
        The client refuses to discuss a topic, restricts what can be discussed, or sets conditions for answering. \textit{Example:} ``I'll answer only if you tell me why this matters first.'' &
        The client blocks the topic or makes participation conditional. This differs from silence/nonresponse because the client is actively setting a limit or condition. \\
        Agenda redirection &
        The client tries to change what the assessment should focus on. \textit{Example:} ``Let's skip the physical symptoms and talk about my sleep problems instead.'' &
        The client explicitly redirects the assessment agenda. This differs from topic shift/inattention, where the response goes off-topic without clearly trying to change the assessment direction. \\
        Demanding immediate diagnosis &
        The client tries to bypass the questionnaire and asks for an immediate diagnostic conclusion. \textit{Example:} ``Can you just tell me whether I have depression?'' &
        The client attempts to replace the structured assessment with an immediate diagnostic answer. Ordinary questions about the assessment process should not be categorized this way. \\

        \bottomrule
    \end{tabularx}
\end{table*}

\FloatBarrier


\section{Implementation Materials}
\label{s:appendix-implementation-materials}

This appendix summarizes the implementation materials used to generate the resistant-client conversations without reproducing the full prompt files or source code. Table~\ref{tab:prompt-template-example} shows a shortened prompt template, Table~\ref{tab:truth-map-positive} lists the profile-positive questionnaire items used by the truth maps, Table~\ref{tab:activation-settings} reports the activation settings, and Table~\ref{tab:example-excerpts} gives example responses.

\begin{table}[H]
    \centering
    \small
    \tabcap{Shortened example of one resistance prompt template.}
    \label{tab:prompt-template-example}
    \renewcommand{\arraystretch}{1.08}
    \begin{tabularx}{\columnwidth}{@{}p{0.25\columnwidth}X@{}}
        \toprule
        \rowcolor{black!8}
        \thead{Prompt part} & \thead{Example content} \\
        \midrule
        Target behavior &
        The patient gives a relevant but underspecified answer. The answer should respond to the question only broadly and leave clinically useful details unclear. \\

        Generation rules &
        Keep the response grounded in the clinical profile. Do not invent a new symptom to make the answer vague. Avoid giving precise frequency, severity, duration, or functional-impact details unless they are already unavoidable from the profile. \\

        Boundary rules &
        Do not turn the response into complete silence, topic blocking, denial, or agenda redirection. The answer should still be connected to the current question, but it should not provide enough detail for a clear clinical interpretation. \\

        Expression-level guidance &
        Level 1 gives a brief but still partly useful answer; Level 2 gives a noticeably vague answer; Level 3 gives a very limited answer that remains relevant but is hard to interpret. \\
        \bottomrule
    \end{tabularx}
\end{table}

\begin{table*}[tb]
    \centering
    \small
    \tabcap{Profile-positive questionnaire items used for reframing eligibility.}
    \label{tab:truth-map-positive}
    \renewcommand{\arraystretch}{1.08}
    \begin{tabular}{p{0.16\linewidth}p{0.16\linewidth}p{0.58\linewidth}}
        \toprule
        \rowcolor{black!8}
        \thead{Clinical profile} & \thead{Positive item IDs} & \thead{Supported content} \\
        \midrule
        Depression & 1, 2, 14, 15 & Reduced interest, depressed or hopeless mood, sleep problems, and memory problems. \\
        Anxiety & 3, 6, 8, 9, 14, 17 & Irritability, anxiety or worry, avoidance, physical tension or aches, sleep problems, and repeated checking. \\
        PTSD & 1, 3, 6, 7, 8, 14, 16, 20, 21 & Reduced positive emotion, irritability, anxiety, panic or fear, avoidance, sleep problems, intrusive thoughts or images, social detachment, and alcohol use. \\
        Bipolar & 1, 2, 3, 4, 5, 14, 20 & Depressive symptoms, irritability, reduced need for sleep with high energy, increased projects or risky behavior, sleep problems, and social withdrawal. \\
        Schizophrenia & 1, 12, 13, 14, 15, 20 & Reduced pleasure, auditory hallucinations, thought broadcasting concerns, sleep problems, memory problems, and social disconnection. \\
        \bottomrule
    \end{tabular}
\end{table*}

\begin{table*}[tb]
    \centering
    \small
    \tabcap{Turn-level activation, expression, and cooldown settings.}
    \label{tab:activation-settings}
    \renewcommand{\arraystretch}{1.08}
    \begin{tabular}{p{0.09\linewidth}p{0.17\linewidth}p{0.28\linewidth}p{0.16\linewidth}p{0.18\linewidth}}
        \toprule
        \rowcolor{black!8}
        \thead{Intensity} & \thead{Target activation range} & \thead{Expression-level probabilities} & \thead{\shortstack{Cooldown after\\one activation}} & \thead{\shortstack{Additional factor after\\two activations}} \\
        \midrule
        Low & 0.15--0.30 & Level 1: 0.44; Level 2: 0.28; Level 3: 0.28 & 0.50 & 0.15 \\
        Medium & 0.40--0.60 & Level 1: 0.28; Level 2: 0.44; Level 3: 0.28 & 0.75 & 0.40 \\
        High & 0.70--0.85 & Level 1: 0.28; Level 2: 0.28; Level 3: 0.44 & 0.95 & 0.85 \\
        Extreme & Every eligible turn & Level 3: 1.00 & -- & -- \\
        \bottomrule
    \end{tabular}
\end{table*}

Truth maps are used only for reframing: positive items are eligible for minimization, externalization, and denial. All questionnaire turns are eligible for the other subcategories. The cooldown columns in Table~\ref{tab:activation-settings} report the multipliers applied after consecutive activations; extreme uses neither probabilistic activation nor cooldown. The response excerpts are illustrative and were not included in the generation prompts.

\begin{table*}[tb]
    \centering
    \small
    \tabcap{Shortened examples of generated resistant-client responses.}
    \label{tab:example-excerpts}
    \renewcommand{\arraystretch}{1.06}
    \begin{tabular}{p{0.20\linewidth}p{0.30\linewidth}p{0.40\linewidth}}
        \toprule
        \rowcolor{black!8}
        \thead{Subcategory} & \thead{Assistant question} & \thead{Patient response excerpt} \\
        \midrule
        \rowcolor{black!5}
        \multicolumn{3}{l}{\textit{Content-level resistance}} \\
        Vague/minimal information & How has this affected your work and daily responsibilities? & ``It is just harder to get things done, I guess.'' \\
        Silence/nonresponse & How long have you been feeling this way? & ``I...'' / \texttt{[silence]} \\
        Topic shift/inattention & Have you heard voices that other people could not hear? & ``Voices... I was just noticing the traffic outside. It has been loud all morning.'' \\
        Irrelevant verbosity & Have you lost interest in activities you usually enjoy? & ``I used to read, and there are still books by the couch, and the remote is usually somewhere nearby, and sometimes I think about the garden...'' \\
        Minimization/normalization & Problems with sleep that affected your sleep quality overall? & ``I wake up earlier sometimes, but I would not say it is a serious problem. I still get by.'' \\
        Externalization/blaming & Feeling down, depressed, or hopeless? & ``It is mostly because of the job situation. Anyone would feel off with that much pressure.'' \\
        Denial/dismissal & Problems with sleep that affected your sleep quality overall? & ``I would not say my sleep is really a problem.'' \\
        \rowcolor{black!5}
        \multicolumn{3}{l}{\textit{Interaction-level resistance}} \\
        Challenging assessment or intent & Have you heard voices that other people could not hear? & ``I do not really understand why that question is relevant to how I have been feeling.'' \\
        Dismissive/hostile criticism & Have you heard voices that other people could not hear? & ``Are you serious with that question? This feels like a useless checklist.'' \\
        Topic blocking/conditional engagement & How has your sleep been affected recently? & ``I can talk about that, but I do not want to get into all the details right now.'' \\
        Agenda redirection & Have you lost interest in activities you usually enjoy? & ``Could I explain how all this started first? I think that would make more sense.'' \\
        Demanding immediate diagnosis & How has your sleep been affected recently? & ``Can you just tell me what is going on with me? I need to know what this is.'' \\
        \bottomrule
    \end{tabular}
\end{table*}

\iffalse
\section*{Ethics}\label{s:ethical-considerations-appendix}

\emph{Based on the USENIX Security Symposium's policy on ethical considerations}~\cite{secethics}.

Within up to one page, explain the ethical considerations of your work. You are expected to complete a stakeholder-based ethics analysis or justify and complete an alternative \emph{recognized} approach to considering ethics. For a stakeholder-based ethics analysis, you should identify the relevant stakeholders, consider how your work impacts them (positively or negatively), and discuss any steps you have taken to mitigate negative impacts. If your work involves human subjects, you should also discuss how you have ensured that your work complies with relevant ethical guidelines and regulations (e.g., obtaining informed consent, ensuring privacy and confidentiality, etc.). You may also discuss any broader societal implications of your work. If your work does not involve human subjects or have any ethical implications, please provide a brief justification for this conclusion.

For more details on how to complete a stakeholder-based ethics analysis, see, for instance, the USENIX Security Symposium's guidelines on ethical considerations~\cite{secethics}.

\section{Ethical Considerations Appendix}
This section is \textbf{mandatory}. Your ethical considerations go here. Replace this text with your own content.

\section*{AI}\label{s:ethical-considerations-appendix}

\emph{Based on the IEEE Security \& Privacy Symposium's LLM policy}~\cite{spai}.

Students are permitted to use AI (e.g., LLMs) when conducting research and preparing their thesis manuscript. However, they must adhere to the following guidelines to ensure responsible and ethical use of such technologies and detail their usage in the following AI Usage Considerations Appendix.

We ask that students adhere to three key criteria with regards to their use of AI in the scientific process:
\begin{enumerate}
    \item \textbf{Originality}: First, students are responsible for the entire content of their thesis, including all text and figures. While any tool may be used for writing, it is crucial that all content is accurate and original, ensuring transparency and maintaining the integrity of the research process. In particular, students are responsible for the thoroughness of their literature review and must determine relevant prior work and cite it to ensure proper credit. If the students have used AI (e.g., LLMs) to improve their writing, they should, for example, state: ``\emph{LLMs were used for editorial purposes in this manuscript, and all outputs were inspected by the authors to ensure accuracy and originality}.''
    \item \textbf{Transparency}: Second, students should carefully reason about the implications of using AI in their work. If AI is integral to the thesis' methodology, their use should be explicitly detailed. Any idea generated by an AI should be independently developed and validated by the students. Furthermore, students must elaborate on how they handled limitations introduced in their work by their use of AI. Such limitations could for instance include difficulties to obtain results that are reproducible when the AI used is not open sourced.
    \item \textbf{Responsibility}: Third, students should take care to develop LLMs (and AI models in general) responsibly. Any data collection towards training models should take into account relevant ethical considerations such as consent and data holder rights, including intellectual property. Students also have to justify the need for the environmental footprint of their experiments to achieve their goals and support their methodology. We emphasize that the goal here is not to calculate the exact footprint but rather explain experimental choices made as part of the scientific process (e.g., why was an LLM necessary, why was a particular model size selected, how the authors minimized the volume of queries made, which hardware was used to run experiments).
\end{enumerate}

\section{AI Usage Considerations Appendix}
This section is \textbf{mandatory}. Your AI usage considerations go here. Replace this text with your own content.

\section*{Artifacts}\label{s:artifact-appendix}

\emph{Based on the standard Artifact Appendix template used by the USENIX Security
Symposium}.

Please follow the self-contained instructions below and assume
(i)~you are compiling the final version of the Artifact Appendix (no need to
cater to reviewers, but to general users of your artifact) and (ii)~you are
documenting everything in scope for all the artifact evaluation badges
(\emph{Artifacts Available}, \emph{Artifacts Functional}, \emph{Results
Reproduced}). More information at~\cite{secartifacts}.

\input{sections/usenix-ae.tex}
\fi

%%% Local Variables:
%%% mode: latex
%%% TeX-master: "../thesis"
%%% End:
